%%%%%%%%%%%%%%%%%%%%%%%%%%%%%%%%%%%%%%%%%%%%%%%%%%%%%%%%%%%
\begin{frame}[fragile]\frametitle{}
\begin{center}
{\Large Demystifying AI}

{\small From Buzzword to Business Enabler}
\end{center}
\end{frame}

%%%%%%%%%%%%%%%%%%%%%%%%%%%%%%%%%%%%%%%%%%%%%%%%%%%%%%%%%%%
\begin{frame}[fragile]\frametitle{What Is AI?, A Business Leader's Definition}

\textbf{Artificial Intelligence:} Machines performing tasks that normally require human intelligence, perceiving, reasoning, deciding, creating.

\vspace{0.5em}
\textbf{The key shift from traditional software:}
\begin{itemize}
  \item Traditional software: A programmer writes every rule explicitly.
  \begin{itemize}
    \item \textit{``If customer buys X, show offer Y''}, a human wrote that rule.
  \end{itemize}
  \item AI / Machine Learning: The system \textit{learns the rules from data automatically}.
  \begin{itemize}
    \item \textit{``We showed 10 million customers their history, the model figured out what to recommend.''}
  \end{itemize}
\end{itemize}

\vspace{0.3em}
\textbf{Plain English:} AI is a junior analyst who has read everything, never sleeps, and gets better every time it gets feedback, but still needs a senior leader to set direction and review critical decisions.

\end{frame}

%%%%%%%%%%%%%%%%%%%%%%%%%%%%%%%%%%%%%%%%%%%%%%%%%%%%%%%%%%%
\begin{frame}[fragile]\frametitle{The Four Types of AI, Business Lens}

    \begin{itemize}
      \item \textbf{Predictive AI:} Learns from historical data to forecast future outcomes.\\
        \textit{Example: Sales forecasting, churn prediction, demand planning.}
      \item \textbf{Generative AI:} Creates new content, text, images, code, audio.\\
        \textit{Example: ChatGPT drafting proposals, DALL-E creating ad visuals.}
      \item \textbf{Agentic AI:} Autonomous systems that plan and execute multi-step tasks.\\
        \textit{Example: AI agent that resolves a customer complaint end-to-end without human input.}
      \item \textbf{Ethical / Responsible AI:} Framework ensuring AI is fair, transparent, and accountable.\\
        \textit{This is not a product, it is a discipline every AI deployment must follow.}
    \end{itemize}

\end{frame}

% %%%%%%%%%%%%%%%%%%%%%%%%%%%%%%%%%%%%%%%%%%%%%%%%%%%%%%%%%%%
% \begin{frame}[fragile]\frametitle{The Four Types of AI, Business Lens}

    % \begin{center}
    % \includegraphics[width=0.8\linewidth,keepaspectratio]{aiaware1}
    % {\tiny (Ref: The Potential of AI to Enhance Business Performance)}
    % \end{center}
 
% \end{frame}


%%%%%%%%%%%%%%%%%%%%%%%%%%%%%%%%%%%%%%%%%%%%%%%%%%%%%%%%%%%
\begin{frame}[fragile]\frametitle{The AI Hierarchy, How the Pieces Fit}

\begin{itemize}
  \item \textbf{AI (broadest):} Any machine mimicking human intelligence
  \item \textbf{Machine Learning (ML):} AI that learns from data, no explicit rules needed
  \item \textbf{Deep Learning:} ML using neural networks; powers image recognition, speech
  \item \textbf{Generative AI:} Deep Learning that \textit{creates} new content
  \item \textbf{LLMs (Large Language Models):} The engine behind ChatGPT, Claude, Gemini, trained on the internet at massive scale
  \item \textbf{AI Agents:} LLMs given tools (web, email, databases) to take autonomous actions
\end{itemize}

\vspace{0.3em}
\textbf{For business leaders:} You don't need to build these. You need to understand which layer solves your problem, and which vendor or team to trust to build it for you.

\end{frame}

%%%%%%%%%%%%%%%%%%%%%%%%%%%%%%%%%%%%%%%%%%%%%%%%%%%%%%%%%%%
\begin{frame}[fragile]\frametitle{From Automation to Intelligence}

    \begin{itemize}
      \item \textbf{Automation:} Rules-based, repetitive task handling. Does exactly what you programmed.
      \item \textbf{AI:} Learns patterns, predicts outcomes, and adapts over time.
      \item \textbf{Business analogy:}
        \begin{itemize}
          \item Automation = ``Doing things right.''
          \item AI = ``Doing the \textit{right} things, smarter.''
        \end{itemize}
      \item \textbf{Example:} Basic chatbot routes calls by keyword (automation). AI chatbot understands intent, sentiment, and context, and improves with every interaction (intelligence).
    \end{itemize}
	
\end{frame}

%%%%%%%%%%%%%%%%%%%%%%%%%%%%%%%%%%%%%%%%%%%%%%%%%%%%%%%%%%%
\begin{frame}[fragile]\frametitle{From Automation to Intelligence}

    \begin{center}
    \includegraphics[width=\linewidth,keepaspectratio]{aiaware2}
    {\tiny (Ref: Automation vs AI Workflows vs AI Agents)}
    \end{center}
 
\end{frame}

%%%%%%%%%%%%%%%%%%%%%%%%%%%%%%%%%%%%%%%%%%%%%%%%%%%%%%%%%%%
\begin{frame}[fragile]\frametitle{Myths vs.\ Reality, What Every Leader Must Know}

\begin{itemize}
  \item \textbf{Myth:} AI understands like a human does.\\
    \textbf{Reality:} AI identifies statistical patterns. It does not ``think.'' It can be confidently wrong.

  \item \textbf{Myth:} AI is always accurate.\\
    \textbf{Reality:} Models ``hallucinate'', generating plausible but incorrect outputs. Human oversight is essential for high-stakes decisions.

  \item \textbf{Myth:} AI will replace our workforce.\\
    \textbf{Reality:} AI replaces repetitive tasks, not roles. Employees who use AI well become 2$\times$--10$\times$ more productive. Demand for skilled people rises.

  \item \textbf{Myth:} We need to train our own AI model.\\
    \textbf{Reality:} For 95\% of business use cases, you \textit{integrate} existing AI (ChatGPT, Claude, Gemini) via APIs on top of your own data.

  \item \textbf{Myth:} Enterprise AI and consumer AI (ChatGPT free) are the same.\\
    \textbf{Reality:} Enterprise AI has data isolation, compliance, SLAs, cost controls, consumer tools do not.
\end{itemize}

\end{frame}

%%%%%%%%%%%%%%%%%%%%%%%%%%%%%%%%%%%%%%%%%%%%%%%%%%%%%%%%%%%
\begin{frame}[fragile]\frametitle{Where AI Is Already in Your Business, Right Now}

Even if your organization hasn't formally adopted AI, it is already present:
\begin{itemize}
  \item \textbf{Email:} Gmail Smart Reply, Outlook Copilot summarizing threads
  \item \textbf{Sales CRM:} Salesforce Einstein, predicting deal close probability
  \item \textbf{Finance:} AI flagging anomalous transactions for fraud detection
  \item \textbf{HR:} Resume screening algorithms, employee attrition prediction
  \item \textbf{Marketing:} Dynamic pricing, personalized ad targeting, content generation
  \item \textbf{Customer Support:} Chatbots handling Tier-1 queries, ticket routing
  \item \textbf{Supply Chain:} Demand forecasting, delivery route optimization
  \item \textbf{Manufacturing:} Predictive maintenance, catching failures before they happen
\end{itemize}

\textbf{The question is no longer ``should we?'', it's ``how fast, and how responsibly?''}

\end{frame}

%%%%%%%%%%%%%%%%%%%%%%%%%%%%%%%%%%%%%%%%%%%%%%%%%%%%%%%%%%%
\begin{frame}[fragile]\frametitle{Enterprise vs.\ Consumer AI, Know the Difference}

\begin{center}
\begin{tabular}{|p{1.5cm}|p{2cm}|p{3cm}|}
\hline
\textbf{Aspect} & \textbf{Consumer (free tools)} & \textbf{Enterprise AI} \\
\hline
Example & ChatGPT free, Gemini free & Azure OpenAI, AWS Bedrock \\
Data Privacy & Shared; may train on your data & Isolated; contractually protected \\
Compliance & None & SOC2, GDPR, ISO 27001 \\
Cost & Free or flat & Token-based; FinOps governance needed \\
Reliability & Best-effort & SLA guaranteed \\
Integration & Manual / copy-paste & API, CRM, ERP, workflows \\
Auditability & None & Full logging, audit trail \\
\hline
\end{tabular}
\end{center}

\vspace{0.3em}
\textbf{Critical for leaders:} Never allow employees to paste confidential data, customer PII, or trade secrets into consumer AI tools. This is a data leakage risk.

\end{frame}

%%%%%%%%%%%%%%%%%%%%%%%%%%%%%%%%%%%%%%%%%%%%%%%%%%%%%%%%%%%
\begin{frame}[fragile]\frametitle{AI What It Cannot Do, Honest Expectations}

Setting realistic expectations prevents failed projects and board-level disappointment:
\begin{itemize}
  \item \textbf{AI cannot reason causally:} It finds correlations, not true cause-and-effect.\\
    \textit{``Ice cream sales and drownings are correlated, the AI doesn't know summer is the cause.''}
  \item \textbf{AI cannot replace human judgment in ethical decisions:} Hiring, promotion, loan approval, legal and moral accountability stays with humans.
  \item \textbf{AI cannot work without data:} Garbage data in, garbage output out. Data quality is a prerequisite.
  \item \textbf{AI cannot explain itself easily:} Most powerful models are ``black boxes.'' Regulated industries (banking, healthcare) need explainability.
  \item \textbf{AI needs maintenance:} Models drift over time as the world changes. Ongoing investment required.
\end{itemize}

\end{frame}

%%%%%%%%%%%%%%%%%%%%%%%%%%%%%%%%%%%%%%%%%%%%%%%%%%%%%%%%%%%
\begin{frame}[fragile]\frametitle{Hands-On Exercise: Prompt Engineering Basics}

    \begin{itemize}
      \item \textbf{Activity:} Explore how AI interprets different prompts.
      \item Compare these two:
        \begin{itemize}
          \item \textit{``Summarize this customer complaint in one line.''}
          \item \textit{``You are a customer success manager. Turn this complaint into a polite, empathetic response email with a resolution offer and follow-up date.''}
        \end{itemize}
      \item Discuss: How does phrasing, role, and context change the output?
      \item \textbf{Key insight:} Better prompts = better results. AI quality is a \textit{leadership skill}.
      \item Reflection: ``Where could clearer prompting improve my team's output this week?''
    \end{itemize}

\end{frame}


%%%%%%%%%%%%%%%%%%%%%%%%%%%%%%%%%%%%%%%%%%%%%%%%%%%%%%%%%%%
\begin{frame}[fragile]\frametitle{Hands-On Exercise: Prompt Engineering Basics}

    \begin{center}
    \includegraphics[width=\linewidth,keepaspectratio]{aiaware3}
    {\tiny (Ref: Prompt Engineering, ai-rockstars.com)}
    \end{center}

\end{frame}

%%%%%%%%%%%%%%%%%%%%%%%%%%%%%%%%%%%%%%%%%%%%%%%%%%%%%%%%%%%
\begin{frame}[fragile]\frametitle{Prompting Patterns Every Leader Should Know}

You do not need to code. You need to communicate with AI effectively.
\begin{itemize}
  \item \textbf{Role Prompting:} ``You are a strategy consultant with 20 years in consumer electronics. Analyze this market data and identify the top 3 growth opportunities.''\\
    \textit{Why it works: Activates the right vocabulary, framing, and level of analysis.}

  \item \textbf{Context + Task:} Always give AI: Who you are, what you need, the format you want, and any constraints.\\
    \textit{``Write a 5-bullet executive summary of this report for a board audience with no AI background.''}

  \item \textbf{Iterate, Don't Accept First Output:} AI output is a first draft, not a final answer. Refine with follow-up prompts.\\
    \textit{``Make it more concise'' / ``Add a risk section'' / ``Change tone to be more assertive.''}

  \item \textbf{Structured Output:} Ask AI for tables, bullet points, or specific sections, output becomes more actionable.
\end{itemize}

\end{frame}


